\documentclass[twoside]{article}
\usepackage{amsmath}
\usepackage{lipsum} % Package to generate dummy text throughout this template
\usepackage[english]{isodate}% http://ctan.org/pkg/isodate
\usepackage{pgfgantt}
\usepackage[show]{ed}
\setlength{\parindent}{0pt}
\usepackage[sc]{mathpazo} % Use the Palatino font
\usepackage[T1]{fontenc} % Use 8-bit encoding that has 256 glyphs
\linespread{1.5} % Line spacing - Palatino needs more space between lines
\usepackage{microtype} % Slightly tweak font spacing for aesthetics
\usepackage{minted}
\usemintedstyle{autumn}
\usepackage[hmarginratio=1:1,top=32mm,left=25mm,right=25mm,columnsep=18pt,bottom=27mm]{geometry} % Document margins
\usepackage{multicol} % Used for the two-column layout of the document
\usepackage[hang, small,labelfont=bf,up,textfont=it,up]{caption} % Custom captions under/above floats in tables or figures
\usepackage{booktabs} % Horizontal rules in tables
\usepackage{float} % Required for tables and figures in the multi-column environment - they need to be placed in specific locations with the [H] (e.g. \begin{table}[H])
\usepackage{hyperref} % For hyperlinks in the PDF

\usepackage{lettrine} % The lettrine is the first enlarged letter at the beginning of the text
\usepackage{paralist} % Used for the compactitem environment which makes bullet points with less space between them
\usepackage{etoolbox}
\patchcmd{\thebibliography}{\section*{\refname}}{}{}{}
\usepackage{abstract} % Allows abstract customization
\renewcommand{\abstractnamefont}{\normalfont\bfseries} % Set the "Abstract" text to bold
\renewcommand{\abstracttextfont}{\normalfont\small\itshape} % Set the abstract itself to small italic text

\usepackage{titlesec} % Allows customization of titles
\renewcommand\thesection{\Roman{section}} % Roman numerals for the sections
\renewcommand\thesubsection{\Roman{subsection}} % Roman numerals for subsections
\titleformat{\section}[block]{\large\scshape\centering}{\thesection.}{1em}{} % Change the look of the section titles
\titleformat{\subsection}[block]{\large}{\thesubsection.}{1em}{} % Change the look of the section titles

\usepackage{fancyhdr} % Headers and footers
\pagestyle{fancy} % All pages have headers and footers
\fancyhead{} % Blank out the default header
\fancyfoot{} % Blank out the default footer
\fancyhead[C]{Naomi Pentrel $\bullet$ \href{mailto:n.pentrel@jacobs-university.de}{n.pentrel@jacobs-university.de} $\bullet$ \today} % Custom header text
\fancyfoot[RO,LE]{\thepage} % Custom footer text
\usepackage{graphicx}
\usepackage{wrapfig}

\usepackage{tocloft}

%\usepackage{csquotes}
%\MakeOuterQuote{"}

\hypersetup{
    bookmarks=true,
    unicode=false,
    pdftoolbar=true,
    pdfmenubar=true,
    pdffitwindow=false,
    pdfstartview={FitH},
    pdftitle={My title},
    pdfauthor={Author},
    pdfsubject={Subject},
    pdfcreator={Creator},
    pdfproducer={Producer},
    pdfkeywords={keyword1} {key2} {key3},
    pdfnewwindow=true,
    colorlinks=false,
    linkcolor=red,
    citecolor=green,
    filecolor=magenta,
    urlcolor=cyan
}

\usepackage{caption}
\usepackage{subcaption}
\usepackage[section]{placeins}
\usepackage{qtree}
\usepackage{gensymb}

\usepackage{setspace}

\usepackage{color}

%\doublespacing
% or:
%\onehalfspacing

%\renewcommand\cftchapafterpnum{\vskip10pt}
%\renewcommand\cftsecafterpnum{\vskip15pt}

%----------------------------------------------------------------------------------------
%	TITLE SECTION
%----------------------------------------------------------------------------------------

\title{\vspace{-15mm}\fontsize{24pt}{10pt}\selectfont\textbf{// All comments are NOT created equal}} % Article title


%----------------------------------------------------------------------------------------
% Macros
%----------------------------------------------------------------------------------------
\newcommand{\sys}{Knowledge Presentation System}

%----------------------------------------------------------------------------------------

\begin{document}
\thispagestyle{empty}
\pagenumbering{roman}
\begin{flushright}
    \includegraphics[scale=1.0]{Logo}
  \end{flushright}
  \vspace{20mm}
  \begin{center}
    \huge
    \textbf{Visualization of Theory Graphs of Semantically Annotated Documents in the Form of Spatial Presentations}
  \end{center}
  \vspace*{4mm}
  \begin{center}
   \Large by
  \end{center}
  \vspace*{4mm}
  \begin{center}
    \Large
    \textbf{Naomi Pentrel}
  \end{center}
  \vspace*{20mm}
  \begin{center}
    \large
    Bachelor Thesis Proposal in Computer Science
  \end{center}
  \vfill
  \begin{flushright}
    \large
    \begin{tabular}{l}
      
      \hline
      Prof. Dr. Michael Kohlhase \\
      \\
    \end{tabular}
  \end{flushright}
  \vspace*{8mm}
  \begin{flushleft}
    \large
    Date of Submission: \today \\
    \rule{\textwidth}{1pt}
  \end{flushleft}
  \begin{center}
    \Large Jacobs University Bremen - School of Engineering and Science
  \end{center}

\newpage
\noindent
  With my signature, I certify that this thesis has been written by me
  using only the indicates resources and materials. Where I have
  presented data and results, the data and results are complete,
  genuine, and have been obtained by me unless otherwise acknowledged;
  where my results derive from computer programs, these computer
  programs have been written by me unless otherwise acknowledged. I
  further confirm that this thesis has not been submitted, either in
  part or as a whole, for any other academic degree at this or another
  institution.

  \vspace{20mm}

    \includegraphics[scale=0.2]{Signature}
 \hfill Bremen, \today
  
\newpage

\thispagestyle{fancy} % All pages have headers and footers

%----------------------------------------------------------------------------------------
%	ARTICLE CONTENTS
%----------------------------------------------------------------------------------------

 \section*{Abstract}
Visualization of knowledge is important to foster learning. To optimize the visualization of knowledge and its interdependencies, we will focus on identifying ways to present knowledge without loosing context. This research aims at taking an existing annotated corpus and presenting its contents in a way that allows the audience to see the dependencies of the covered topics. This form of presenting the knowledge should make it easy to examine materials that form the basis for the topics the audience is interested in.\\

Consider the typical lecture one attends as a student. Towards the end of the lecture students might have difficulties remembering earlier topics. When introducing a new topic it would be ideal to have a simple way to find the knowledge dependencies and present students with a simple way to catch up.\\

The OMDoc (Open Mathematical Documents) format is a content markup scheme for mathematical documents.\cite{OMDoc book} Using OMDoc, this research aims to visualize the connections of some of what is written about Mathematics or other fields in a way that allows students to learn the concepts the current topic depends on should they require to refresh their memories. As such it will foster learning and could transform how we interact with course material as a whole. It is easy to see that this research would also have many other real world applications anywhere where knowledge is supposed to be transferred.\\ 

  \newpage
  \tableofcontents

  \clearpage
  \pagenumbering{arabic}

  \section{Introduction}

The Semantic Web approach aims to transform the World Wide Web (WWW) from a mere information space to an information web that is processable by machines.\cite{SemWeb} There exists an interesting similarity in other fields such as that of new presentation tools that aim to enable you to present information in a more interconnected way that provides context. Both developments strive to connect pieces of information and through providing context elevate pure information that is provided to knowledge.\\

These two developments have more in common than one would think. Looking at the history of the WWW, the era of Web 2.0, where everyone is now able to generate content lead to a lot of information but comparatively little knowledge.\cite{Weller} Here the processing of information to provide knowledge is lacking. Similarly I would argue that the existance of tools like Powerpoint have lead to virtually everyone doing presentations. As everyone who has witnessed many presentation knows, these are easily compared to blogs: they contain some knowledge but are often badly delivered and lack structure and context.\\

Semantic Web is aiming to change the former, whereas new, more visual presentation tools provide a solution to help provide presentations that are easier to process and thus will facilitate knowledge transfer. This research intends to bridge the gap and provide a tool for contextual visualization of information. It adapts some of the concepts of the semantic web on a smaller scale.

In this paper we want to go from merely displaying information to visualizing information to provide knowledge by displaying the network/context a piece of information is in. In this we will not try to semantically represent the meaning of the information that is shown but rather on showing the network, i.e. context the information is in.

\ednote{INSERT GRAPH that leads from info to knowledge
\cite{Probst} \\}

"Knowledge is the information necessary to support intelligent reasoning." \cite{Kohlhase} That said, it is imperative to facilitate knowledge transfer.  Thus, we will examine how visually representing information and its context can make knowledge transfer easier. For this we will be examining the OMDoc (Open Mathematical Documents) format as the example that will provide us with an annotated set of documents of which we will visualize the dependencies of the contained information. We will however not visualize the dependencies as a simple graph but use an API such as Impress.js to allow for us to show the current piece of information of interest in context with dependent information.\\

To exemplify this: Annie, a young student learning Spanish is watching her teacher give a presentation on how to form sentences. She wants to form the sentence 'I read books'. After struggling and looking up the word for reading she forms the sentence 'Yo leer libros'. For this Annie had to find the dependency on knowing the vocabulary words for 'I', 'read', and 'books'. However Annie made a mistake and did not know that the verb 'leer' needs to be conjugated appropriately. Her teacher tells her this. Now Annie has to try to find the information on how to conjugate verbs. This shows the dependencies of different pieces of information, i.e. the formulation of a sentence depends on knowledge of grammar and vocabulary. In an ideal presentation this information should have been linked to the current piece of information so that Annie could have directly looked it up instead of having to search for it.

After reading this contrived story, I leave the reader to think about examples in his or her own life where he or she would have appreciated having the context of the information he or she was looking at at the time. Many of these stories exist in real life and it is therefore important to provide this context to facilitate learning and knowledge transfer as a whole. 


\ednote{-addresses fellow experts from your field (but not from your particular topic of research). }
\ednote{- Here you should technically connect to the main concepts from that field and give an outline of your project, stating the research/engineering question that you want to get answered by your project.}
\ednote{Why is this important for the world? How will this improve learning? ->More effective education. }
\ednote{(target size: 1 page)}
\ednote{http://rhizomik.net/html/~roberto/thesis/html/KnowledgeRepresentation.html}
\ednote{link}
% http://reader.eblib.com/(S(gkwt4h53aadjotp03mi25zsi))/Reader.aspx?p=655998&o=1706&u=tkyUhRRU8oBFV92xtGUdWNHfafo%3d&t=1416427415&h=AAEABAAD4041366D9CC8528AF618690B20349073&s=31025691&ut=5736&pg=1&r=img&c=-1&pat=n&cms=-1&sd=2

\section{Preliminaries}
To understand the problem of information that lacks context, this part will provide the necessary background information to allow us to derive a possible solution. While doing so it will demonstrate why this research is important and how it will improve the transfer of knowledge.

\ednote{ This part should make clear which question, exactly, you are pursuing, and why your project is relevant/interesting. This is the place to cite relevant literature. Where does your project extend the state of the art? what weaknesses in known approaches to you hope to overcome? If you have carried out preliminary experiments, describe them here.}

\ednote{(target size: 3-4 pages)}
\ednote {explain primitives and apply to prezi.}

\subsection{OMDoc}
(half a page)
\vspace*{8cm}
\subsection{Status of Information within Discourse}

In linguistics the concept of information packaging is know. Within this, one discriminates between \textit{familiar} and \textit{old} information. Familiar information is shared by speaker and addressee whereas new information not. In addition one distinguishes information that is old or new with respect to the discourse or with respect to the addressee.\\

[1] \textit{"My sister went to the circus the other day; \underline{she} said \underline{it} was brilliant."}\\

In this example in the first part of the sentence information pertaining to my sister and to a circus are introduced, i.e. \textit{discourse-new}. In the second part the underlined parts are considered \textit{discourse-old} since they have already been introduced.\cite{Cambridge}\\

These concepts can be adapted to the situation of teaching Mathematics or Computer Science to students in a class. In general the information that the speaker is sharing with the addressees/students is \textit{discourse-new}\ednote{call course-new?} and depend on  information that is \textit{discourse-old}. I would like to introduce the term \textit{discourse-ancient} which covers the situation where the addressee has difficulties following the speaker since the \textit{discourse-new} information depends on \textit{discourse-old} information that might be 'too old', i.e. ancient, to be easily remembered.\\

\subsection{Primitives}

\subsection{Context}

\ednote{prezi - cell example.}

\subsection{Spatial Narrative}
The development of of presentation tools like Impress.js and Prezi, "a virtual whiteboard that transforms presentations from monologues into conversations: enabling people to see, understand, and remember ideas"\cite{Prezi}, has  changed how people think about presentations. When doing a good prezi presentation one has to understand the topic one is presenting and think about how one can portray connections between content visually.\\

The concept of Spatial Narrative stems from the establishment of frameworks for "the creation of computer-assisted flexibile 'guided tours' based on the thematically and narrative linking of a set of locations within an area into a 'spatial narrative'". An example for this is the experience of a tour guide for a certain topic.\cite{Narrative}\\

The benefits of spatial narrative is that the audience gains a more thorough understanding of the subject and has a better understanding of the interconnections between different pieces of information. By providing the audience with a story as a narrative one also makes use of the concept of Storytelling which is largely accredited with the benefit that the audience will have an easier time following the presenter. That is because our brains are not made to memorize a lot of unconnected pieces of information; it is far easier for us to remember information if it comes in story form. \ednote{find source}\\

In this research we will use this knowledge to form an interconnected web of information that tells a visual story to facilitate the transfer of knowledge.

\subsection{Different Representations of Data}
What's the best representation one could do?
(1 page)


\section{Planned Investigation}

  This is the technical core of the proposal. Here you lay out your
  plans of how you want to answer your research question specify your
  design of experiments or simulations, point out difficulties that
  you expect to encounter, etc.

How will this be implemented?
Using which frameworks?
  

\subsection{Forming Visualizations out of Knowledge Graphs which Convey Meaning}
(1 page)
\newpage

\subsection{Ordered Knowledge Graphs}
(1 page)
\newpage

\subsection{Meaningful Representation of Knowledge Graphs}
(1 page)

\ednote{evaluation criteria from lecture slides}
\newpage

\section{Work Plan}
\subsection{WP1: Understand API of Impress.js}

The first step to implementing a \sys is to understand the tool we want to represent the knowledge in later. Since Prezi at the current time does not have an API that would allow for the generation of a whole presentation without the editor, this research will employ another tool that is inspired by the idea behind Prezi. Bartek Szopka's \textbf{Impress.js} is an open-source presentation framework which is based on CSS3 transforms and transitions.\cite{Impress} \ednote{Impressionist and jimpress}\\

\subsection{WP2: Using OMDoc to create ordered information graphs}
The idea is to use the General Computer Science Lecture notes as an experiment. The General Computer Science Lecture notes are an annotated set of documents. Using OMDoc, we will create an information graph out of the annotated set of documents. This information graph should be ordered based on dependencies of knowledge.\\

\subsection{WP3: Combining OMDoc and Spatial Narration}
Once we have an ordered information graph we will visualize the information in different ways and analyze each way. The outcome of the analysis should be a model which we will then use to visualize the information in WP4.\\

\subsection{WP3: Implement OMDoc examples in Impress.js}
After the best model for spatial narration is found, we will implement a system that will take the ordered information graph and transform it into a spatial presentation using Impress.js.

\subsection{WP4: Evaluation}
Once we have results, a qualitative analysis of the visualization produced by the implemented \sys should be perfomed. In addition we will determine whether the implemented \sys meets the goals outlined in the Evaluation Criteria. Lastly we will measure the speed of the system.

\section{Evaluation Criteria}

\ednote{When do I feel successful? when I have running system (or I can tell why it didn't/couldn't work) or when users can understand it and reconstruct the original graph. }

The evaluation criteria for the research at hand are as follows. 
\begin{enumerate}
\item Does the \sys that was developed allow for the creation of a full presentation out of a given annotated document?
\item  How fast does the \sys work? Can this be used in real time? Could a student request a presentation on a certain topic and could this be generated on the fly within seconds?
\item Does a qualitative analysis yield the result that the developed system provide a better system for knowledge transfer?
\item What are the possitive effects on the combination of Spatial Narrative and OMDoc?
\item Are there negative side effects?
\end{enumerate}
\ednote{I am missing SMART goals... Can I use a scale from 0 - 10 for each point?}


\ednote{(target size: 1-2 page)}
\newpage
  \section{Timeline}

\textbf{Milestones}
\begin{itemize}
\item \textbf{Until February 15th}: Complete preliminary development to create a presentation from an XML
\item \textbf{Until April 1st}: Complete development of organized presentations that are generated from ordered knowledge graphs.
\item \textbf{Until May 1st}: Testing and Performance Measurement.
\item \textbf{Until May 10th}: Complete development of organized presentations that are generated from ordered knowledge graphs.
\end{itemize}

\ednote {wp 4 - thesis writing}

\begin{figure}[h]
\begin{center}
\begin{ganttchart}[y unit title=0.4cm,
y unit chart=0.5cm,
vgrid,hgrid, 
title label anchor/.style={below=-1.6ex},
title left shift=.05,
title right shift=-.05,
title height=1,
bar/.style={fill=gray!50},
incomplete/.style={fill=white},
progress label text={},
bar height=0.7,
group right shift=0,
group top shift=.6,
group height=.3,
group peaks={}{}{.2}]{24}
%labels
\gantttitle{Proposed Timeline}{24} \\
\gantttitle{December}{4} 
\gantttitle{January}{4} 
\gantttitle{February}{4} 
\gantttitle{March}{4} 
\gantttitle{April}{4} 
\gantttitle{May}{4} \\
%tasks
\ganttbar{WP 1}{2}{10} \\
\ganttbar{WP 2}{11}{16} \\
\ganttbar{WP 3}{17}{20} \\
\ganttbar{WP 4}{21}{21.5}

%relations 
\ganttlink{elem0}{elem1} 
\ganttlink{elem1}{elem2} 
\ganttlink{elem2}{elem3} 
\end{ganttchart}
\end{center}
\caption{Gantt Chart}
\end{figure}

  \section{Conclusions}

  Summarize the main aspects of the proposal.

  (target size: 1/2 page)

\newpage
\section{References}
\begin{thebibliography}{} % Bibliography - this is intentionally simple in this template

\bibitem {SemWeb}
Berners-Lee, Tim. "Semantic Web Road Map." Semantic Web Roadmap. World Wide Web Consortium, Sept. 1998. Web. 23 Nov. 2014.
\bibitem {Prezi}
"Hi, We're Prezi." Prezi.com. N.p., n.d. Web. 23 Nov. 2014.
\bibitem {Cambridge}
Huddleston, Rodney, Geoffrey K. Pullum, Laurie Bauer, Betty J. Birner, and Ted Briscoe. "Chapter 16 Information Packaging." The Cambridge Grammar of the English Language. Cambridge: Cambridge UP, 2002. N. pag. Print.
\bibitem {Kohlhase}
Kohlhase, Michael. "320441 CompLog Lecture Notes." (n.d.): n. pag. 17 Nov. 2014. Web. 19 Nov. 2014.
\bibitem {OMDoc book}
Kohlhase, Michael, and Alan Bundy. OMDoc - an Open Markup Format for Mathematical Documents. Berlin: SpringerLink, 2006. Print.
\bibitem{Narrative}
M{\o}ller-Jensen, Lasse. "Mobile GIS Strategies for Disseminating Thematic Tourist Information: Examples of Spatial Narratives." Diss. U of Copenhagen, n.d. Web. 24 Nov. 2014. \href{http://ifgi.uni-muenster.de/archives/tipugg/paper/tipugg2008_Lasse.pdf}{http://ifgi.uni-muenster.de/archives/tipugg/paper/tipugg2008\_Lasse.pdf}.
\bibitem {Probst}
Probst, Gilbert, Steffen Raub, and Kai Romhardt. \newblock Wissen Managen. \newblock Wiesbaden: S.n., 1997. Print.
\bibitem {Impress}
Szopka, Bartek. "Impress.js." Impress.js. Github, 28 Dec. 2011. Web. 23 Nov. 2014.
\bibitem {Weller}
Weller, Katrin. Knowledge Representation in the Social Semantic Web. New York: De Gruyter Saur, 2010. Print.
\end{thebibliography}

\ednote{Research field: information systems}

\end{document}